\label{ch:data-release}

\section{Publications}

\begin{refbox}
\begin{itemize}
  \item \textbf{Data release.} Daya Bay Collaboration, \emph{Full Data
    Release of the Daya Bay Reactor Neutrino Experiment}, v1.0.0, Zenodo,
    DOI:10.5281/zenodo.17587229 (2025) \parencite{DayaBayZenodo2025}. The
    analysis-dataset (TSV) tier is fetched from the official GitHub mirror
    \code{github.com/dayabay-experiment/dayabay-data-official} by
    \code{notebooks/external/reactor/dayabay/DayaBay1\_generator.ipynb}
    (\code{zenodo.org} itself returns HTTP 403 from this project's
    development sandbox) into \code{data/detector/dayabay/}, in the tidy
    \code{.csv} form \texttt{tpeanuts.detector.dayabay} consumes.
  \item \textbf{Main precision measurement.} F.~P.~An et al.\ (Daya Bay
    Collaboration), \emph{Precision Measurement of Reactor Antineutrino
    Oscillation at Kilometer-Scale Baselines by Daya Bay}, Phys.\ Rev.\
    Lett.\ \textbf{130}, 161802 (2023) \parencite{An2023DayaBayPrecision}.
    Publishes the $\sin^2 2\theta_{13}$/$\Delta m^2_{ee}$ best fit compared
    against this document's own inference in \cref{ch:results}.
  \item \textbf{Discovery paper.} F.~P.~An et al.\ (Daya Bay
    Collaboration), \emph{Observation of Electron-Antineutrino
    Disappearance at Daya Bay}, Phys.\ Rev.\ Lett.\ \textbf{108}, 171803
    (2012) \parencite{An2012DayaBayDiscovery}. Source of the original
    near/far rate-comparison strategy \cref{ch:inference-process}
    reproduces.
  \item \textbf{Flux and spectrum measurement.} F.~P.~An et al.\ (Daya Bay
    Collaboration), \emph{Improved measurement of the reactor antineutrino
    flux and spectrum at Daya Bay}, Chin.\ Phys.\ C \textbf{41}, 013002
    (2017) \parencite{An2017DayaBayFluxSpectrum}. Source of an
    independently published per-detector IBD rate table (its Table 5) used
    in \cref{ch:results} as an external cross-check of this project's
    absolute-normalization audit, and of the Reactor Antineutrino Anomaly
    measurement discussed in \cref{ch:detector-implementation}.
  \item \textbf{External comparison: NuFIT.} I.~Esteban et al., \emph{The
    fate of hints: updated global analysis of three-flavor neutrino
    oscillations}, JHEP \textbf{12}, 216 (2024), arXiv:2410.05380
    \parencite{EstebanNuFit2024} -- base analysis publication for NuFIT 6.1
    (2025), \url{www.nu-fit.org}. Source of the fixed solar parameters used
    throughout \cref{ch:detector-implementation}, and of the independent
    NuFIT $\theta_{13}$ comparison in \cref{ch:results}.
\end{itemize}
\end{refbox}

\section{Data-release structure and format conventions}

\begin{theorybox}
The official release ships four parallel data-format tiers (\code{npz},
\code{hdf5}, \code{tsv.bz2}, \code{root}), each containing the identical
physical content; \texttt{tpeanuts} consumes the \code{tsv.bz2} tier only,
since it requires no format-specific reader beyond the standard-library
\code{bz2} module and \code{pandas}. Three real naming conventions appear
in the tables catalogued below:
\begin{itemize}
  \item \code{\{filename\}\_\{key\}.tsv.bz2} -- a single array, named by
    its key (e.g.\ \code{detector\_iav\_matrix\_iav\_matrix.tsv.bz2}).
  \item \code{\{filename\}.tsv/\{key\}.tsv.bz2} -- a folder of related
    arrays, one file per key (e.g.\ per-detector IBD spectra).
  \item \code{\{filename\}.tsv/\{filename\}\_\{key\}.tsv.bz2} -- the same,
    with the folder name repeated in each file for extra clarity when
    nested paths are otherwise ambiguous.
\end{itemize}
Real parameter files (\code{parameters/*.yaml}) additionally carry a
\code{format} field describing their own uncertainty convention
(\code{value}: point value only; \code{sigma\_absolute}: value $\pm$
absolute 1$\sigma$; \code{sigma\_percent}: value $\pm$ percent 1$\sigma$)
and a \code{state} field (\code{fixed} or \code{variable}, i.e.\ whether the
official combined analysis treats that parameter as a free nuisance).
\end{theorybox}

\section{Table catalogue}

\Cref{tab:dayabay-tables-geometry,tab:dayabay-tables-response,tab:dayabay-tables-observed,tab:dayabay-tables-truth}
list every real table this project fetches and the \texttt{tpeanuts}
loader that reads it (all in \texttt{tpeanuts.detector.dayabay.io}),
grouped by category; the physics each table encodes is developed in full
in \cref{ch:detector-implementation}.

\newcommand{\dbcell}[1]{\parbox[t]{0.47\textwidth}{\raggedright #1}}

\begin{table}[htb]
\centering
\caption{Real geometry and reactor-flux tables.}
\label{tab:dayabay-tables-geometry}
\begin{tabular}{@{}l l l@{}}
\toprule
\textbf{Cached file} & \textbf{Real contents} & \textbf{Loader} \\
\midrule
\code{baselines.csv} & \dbcell{(detector, reactor) baseline in km; real $8\times6=48$-row geometry matrix.} & \pyfunc{load\_baselines} \\[4pt]
\code{reactor\_parameters.csv} & \dbcell{Real fission fractions, energy released per fission (4 isotopes), nominal thermal power $2.895$~GW.} & \pyfunc{load\_reactor\_parameters} \\[4pt]
\code{huber\_mueller\_spectra.csv} & \dbcell{Real per-isotope antineutrino spectra \parencite{Huber2011,Huber2012Erratum,Mueller2011}, $50$~keV mesh, $1.75$--$13.0$~MeV.} & \pyfunc{load\_huber\_mueller\_spectra} \\[4pt]
\code{nonequilibrium\_correction.csv} & \dbcell{Real per-isotope non-equilibrium flux correction ($^{235}$U, $^{239}$Pu, $^{241}$Pu only).} & \pyfunc{load\_nonequilibrium\_correction} \\[4pt]
\code{snf\_correction.csv} & \dbcell{Real per-reactor spent-nuclear-fuel flux correction, all 6 cores.} & \pyfunc{load\_snf\_correction} \\[4pt]
\code{neutrino\_rate\_weekly.csv} & \dbcell{Real weekly-resolved per-reactor antineutrino rate, full 2011--2020 data-taking period.} & \pyfunc{load\_neutrino\_rate\_weekly} \\
\bottomrule
\end{tabular}
\end{table}

\begin{table}[htb]
\centering
\caption{Real detector-response and IBD-cross-section tables.}
\label{tab:dayabay-tables-response}
\begin{tabular}{@{}l l l@{}}
\toprule
\textbf{Cached file} & \textbf{Real contents} & \textbf{Loader} \\
\midrule
\code{n\_protons.csv} & \dbcell{Real target free-proton count per detector ($\approx 1.43\times10^{30}$ nominal, $\pm$ small per-detector correction).} & \pyfunc{load\_n\_protons} \\[4pt]
\code{eres\_parameters.csv} & \dbcell{Real 3-term Gaussian energy-resolution coefficients $(a,b,c)$.} & \pyfunc{load\_eres\_parameters} \\[4pt]
\code{detector\_efficiency.txt} & \dbcell{Real flat IBD-selection efficiency, $\varepsilon\approx0.8025$.} & \pyfunc{load\_detector\_efficiency} \\[4pt]
\code{ibd\_constants.csv} & \dbcell{Real IBD cross-section constants: $f=1.0$, $g=1.2701$, $f_2=3.706$, PhaseSpaceFactor$=1.71465$.} & \pyfunc{load\_ibd\_constants} \\[4pt]
\code{iav\_matrix.csv} & \dbcell{Real $240\times240$ IAV energy-redistribution matrix, 0.05~MeV bins, 0--12~MeV, columns summing to 1.} & \pyfunc{load\_iav\_matrix} \\[4pt]
\code{lsnl\_curve\_nominal.csv} & \dbcell{Real LSNL nominal nonlinearity curve $f(E)=E_{\rm reco}/E_{\rm true}$.} & \pyfunc{load\_lsnl\_curve} \\[4pt]
\code{lsnl\_curve\_pulls.csv} & \dbcell{Real 4 LSNL systematic pull curves $f_k(E)$, $k=0..3$.} & \pyfunc{load\_lsnl\_curve\_pulls} \\
\bottomrule
\end{tabular}
\end{table}

\begin{table}[htb]
\centering
\caption{Real observed-data and background tables (8AD period).}
\label{tab:dayabay-tables-observed}
\begin{tabular}{@{}l l l@{}}
\toprule
\textbf{Cached file} & \textbf{Real contents} & \textbf{Loader} \\
\midrule
\code{ibd\_spectra/ibd\_spectrum\_\{AD\}.csv} & \dbcell{Real observed IBD prompt-energy candidate spectra, 0.05~MeV bins, 0--12~MeV, one file per detector.} & \pyfunc{load\_ibd\_spectrum} \\[4pt]
\code{background\_rates\_8AD.csv} & \dbcell{Real background rates + uncertainties, 5 categories $\times$ 8 detectors.} & \pyfunc{load\_background\_rates} \\[4pt]
\code{background\_shapes/spectrum\_shape\_\{cat\}\_\{AD\}.csv} & \dbcell{Real normalised background spectral shapes (integrate to 1), same binning as the IBD spectra.} & \pyfunc{load\_background\_shape} \\[4pt]
\code{exposure\_8AD.csv} & \dbcell{Real per-detector effective livetime (\code{livetime}$\times$\code{eff}), summed over 8AD-period days.} & \pyfunc{load\_exposure} \\[4pt]
\code{final\_erec\_bin\_edges.csv} & \dbcell{Real (non-uniform) analysis binning, 0.7--12.0~MeV, 27 edges.} & \pyfunc{load\_final\_erec\_bin\_edges} \\
\bottomrule
\end{tabular}
\end{table}

\begin{table}[htb]
\centering
\caption{Real published-truth and normalization tables.}
\label{tab:dayabay-tables-truth}
\begin{tabular}{@{}l l l@{}}
\toprule
\textbf{Cached file} & \textbf{Real contents} & \textbf{Loader} \\
\midrule
\code{survival\_probability\_truth.csv} & \dbcell{Daya Bay's own real published best fit ($\sin^2 2\theta_{13}$, $\Delta m^2_{32}$) and solar input ($\sin^2 2\theta_{12}$, $\Delta m^2_{21}$), for comparison only.} & \pyfunc{load\_survival\_probability\_truth} \\[4pt]
\code{global\_normalization.txt} & \dbcell{Daya Bay's own real nominal free-normalization value (1.0).} & \pyfunc{load\_global\_normalization} \\
\bottomrule
\end{tabular}
\end{table}

\begin{notebox}
Not fetched/reproduced: the 6AD/7AD sub-periods (only 8AD); the
LSNL/IAV/background official \emph{uncertainty} parameters as a full
correlated covariance (their real nominal/central-value curves and rates
\emph{are} used, see \cref{tab:dayabay-tables-response}); and per-detector/per-hall
correlation structure for the background-rate systematics (a single
correlated scale per category is used instead, \cref{ch:inference-process}).
\end{notebox}
